\documentclass[10pt,aspectratio=169]{beamer}

\usepackage{spbu}

\usepackage[T2A]{fontenc}
\usepackage[utf8]{inputenc}
\usepackage[english,russian]{babel}

\title{Шаблон презентации для СПбГУ}
\subtitle{\LaTeX/Beamer}
\author{Фамилия Имя Отчество}


\begin{document}

    \maketitle

    \section{Введение}

    \begin{frame}{Пакет Beamer для подготовки презентации}
        \begin{itemize}
            \item Предполагается, что вы умеете пользоваться \LaTeX, если это не так, вы можете начать
                \href{https://www.latex-project.org/}{ознакомление с ним здесь}.
            \item Beamer --- один из самых популярных и мощных пакетов для подготовки презентаций в \LaTeX.
            \item У пакета Beamer есть подробное
                \href{http://www.ctan.org/tex-archive/macros/latex/contrib/beamer/doc/beameruserguide.pdf}{руководство пользователя}.
            \item Данная презентация представляет из себя шаблон/пример, подготовленный с использованием пакета Beamer с фирменной символикой и цветовой гаммой СПбГУ.
        \end{itemize}
    \end{frame}

    \section{Редактирование}

    \begin{frame}[fragile]{Использование Beamer}
        Для создания простейшего документа с использованием пакета Beamer достаточно воспользоваться следующим кодом:
        \begin{block}{Минимальный документ}
\begin{lstlisting}
\documentclass{beamer}
\begin{document}
    \begin{frame}{Hello, world!}
    \end {frame}
\end{document}
\end{lstlisting}
        \end{block}
    \end{frame}
    
    \begin{frame}[fragile]{Использование Beamer}
   		Для использования данного шаблона необходимо подключить стилевой файл \verb|spbu.sty|, который идет в комплекте:
   		\begin{block}{Использование представленного шаблона}
\begin{lstlisting}
\documentclass[10pt,aspectratio=169]{beamer}

\usepackage{spbu}

\usepackage[T2A]{fontenc}            % Внутренняя кодировка T2A
\usepackage[utf8]{inputenc}          % Указывает, что main.tex в кодировке UTF-8
\usepackage[english,russian]{babel}  % Подключает поддержку многоязычного текста

\begin{document}
	\maketitle
	\begin{frame}{Hello, world!}
	\end {frame}
\end{document}
\end{lstlisting}
		\end{block}
   	\end{frame}

    \begin{frame}[fragile]{Титульный слайд}
        Для того, чтобы отобразить информацию на титульном слайде, необходимо вызвать несколько команд в преамбуле:
        \begin{block}{Команды для вывода информации на титульный слайд}
\begin{lstlisting}
\title{Заголовок}
\subtitle{Подзаголовок}
\author{Первый автор, Второй автор}
\date{Дата, по умолчанию - сегодня}
\end{lstlisting}
        \end{block}
        Затем, после команды \verb|\begin{document}|, необходимо добавить команду \verb|\maketitle|, как было показано на предыдущем слайде.
    \end{frame}

    \begin{frame}[fragile]{Создание слайда с информацией}
        \framesubtitle{Руководство}
        Для создания простого слайда, необходимо воспользоваться следующим кодом:
        \begin{block}{Код для создания слайда}
            \begin{lstlisting}[language=TeX]
\begin{frame}
    \frametitle{Пример слайда}
    \framesubtitle{Результат}
    Содержание слайда...
\end{frame}\end{lstlisting}
        \end{block}
        Результат представлен на следующем слайде.
    \end{frame}

    \begin{frame}
        \frametitle{Пример слайда}
        \framesubtitle{Результат}
        Содержание слайда...
    \end{frame}

    \begin{frame}[fragile]{Добавление изображений}
        \begin{columns}
            \begin{column}{0.7\textwidth}
                Добавление изображений работает как в стандартном \LaTeX:
                \begin{block}{Добавление изображения}
\begin{lstlisting}
\usepackage{graphicx}
% ...
\includegraphics
[width=\textwidth]{images/default}
\end{lstlisting}
                \end{block}
            \end{column}
            \begin{column}{0.3\textwidth}
                \includegraphics[width=\textwidth]{source/logo-name}\\
            \end{column}
        \end{columns}
    \end{frame}
    
    \begin{frame}[fragile]{Добавление изображений}
    	\framesubtitle{Подписи, метки, перекрестные ссылки для изображений}
    	\begin{columns}
    		\begin{column}{0.7\textwidth}
    			Используем окружение \verb|figure|, после этого можно добавить подпись для изображения, используя команду \verb|\caption|, и метку, используя команду \verb|\label|. 
    			\begin{block}{Добавление изображения}
\begin{lstlisting}
\begin{figure}[h!]
	\centering
	\includegraphics[width=\textwidth]{source/logo-name}
	\caption{Подпись}
	\label{fig:logo}
\end{figure}
\end{lstlisting}
    			\end{block}
    		\end{column}
    		\begin{column}{0.3\textwidth}
    			\begin{figure}[h!]
    				\centering
    				\includegraphics[width=\textwidth]{source/logo-name}
    				\caption{Подпись}
    				\label{fig:logo}
    			\end{figure}
    		\end{column}
    	\end{columns}
		Теперь можно сослаться на рис.~\ref{fig:logo}, используя команду \verb|\ref{fig:logo}|.
    \end{frame}

    \begin{frame}[fragile]{Добавление формул}
        Формулы добавляются также как и в стандартном \LaTeX:
        \begin{itemize}
            \item Строчные формулы добавляются с использованием символа доллар (\verb|$|), например\\
            \verb|$\Delta =\nabla^{2}$|: $\Delta =\nabla^{2}$.
            \item Вынесенные формулы добавляются, например, с помощью двух символов доллара, либо с использованием окружения
            \verb|equation|, либо любым другим предпочтительным для вас способом.
        \end{itemize}
    \end{frame}

    \begin{frame}[fragile]{Добавление формул}
        \framesubtitle{Пример вынесенной формулы}

        \begin{block}{Добавление вынесенной формулы}
\begin{lstlisting}
\begin{equation}
    \frac{\partial u}{\partial t}
    -a^2\left(\frac{\partial^2 u}{\partial x^2_1}
    +\frac{\partial^2 u}{\partial x^2_2}
    +\ldots+\frac{\partial^2 u}{\partial x^2_n}\right)=f\left(x,t\right).
    \label{eq:eq1}
\end{equation}
\end{lstlisting}
        \end{block}

        Результат:
        \begin{equation}
            \frac{\partial u}{\partial t}
            -a^2\left(\frac{\partial^2 u}{\partial x^2_1}
            +\frac{\partial^2 u}{\partial x^2_2}
            +\ldots+\frac{\partial^2 u}{\partial x^2_n}\right)=f\left(x,t\right).
            \label{eq:eq1}
        \end{equation}
        Ссылка на формулу (\ref{eq:eq1}).
    \end{frame}

    \begin{frame}[fragile]{Разделение на столбцы}
        Разделение информации на столбцы осуществляется следующим способом:
        \begin{block}{Разделение информации на столбцы}
\begin{lstlisting}
\begin{columns}
    \begin{column}{0.6\textwidth}
        первый столбец,
    \end{column}
    \begin{column}{0.3\textwidth}
        второй столбец.
    \end{column}
    % Может быть и больше!
\end{columns}
\end{lstlisting}
        \end{block}

        Далее представлен результат:
        \begin{columns}
            \begin{column}{0.6\textwidth}
                первый столбец,
            \end{column}
            \begin{column}{0.3\textwidth}
                второй столбец.
            \end{column}
        \end{columns}
    \end{frame}

    \begin{frame}[fragile]
        \frametitle{Шрифты}
        Для представления различной информации, уместно использовать соответствующие ей шрифты:
        \begin{itemize}
            \item {\verb|\textrm{}|: \textrm{стандартный шрифт;}}
            \item {\verb|\textsf{}|: \textsf{шрифт семейства sans-serif;}}
            \item {\verb|\texttt{}|: \texttt{моноширинный шрифт.}}
        \end{itemize}
    \end{frame}
    
    \begin{frame}[fragile]
    	\frametitle{Листинг}
    	
    	Для вставки программного кода может использоваться окружение \texttt{lstlisting}. Пример:
\begin{lstlisting}[language=TeX]
\begin{lstlisting}
	# Это пример кода на Python
	def hello_world():
		print("Hello, World!")
\ end{lstlisting}                % Тут добавлен лишний пробел после символа "\"
\end{lstlisting}
	   	Результат:
\begin{lstlisting}
# Это пример кода на Python
def hello_world():
	print("Hello, World!")
\end{lstlisting}

Для вставки строчных примеров кода, можно использовать команду\linebreak \texttt{\textbackslash verb|некоторый код...|}, в результате получим: \verb|некоторый код...|
    \end{frame}
    
    \section{Заключение}

    \begin{frame}
        \frametitle{Заключение}
        \begin{itemize}
            \item Автор будет рад вашему вкладу в репозиторий на \href{github.com}{GitHub}.
            \item Вопросы, пожелания, предложения можно отправлять на
            \href{mailto:gleb.alcybeev@gmail.com}{электронную почту} или писать в \href{https://github.com/gleb-kun/SPbU-BeamerSlides/issues}{обсуждениях на GitHub}.
        \end{itemize}
    \end{frame}

    \backmatter

\end{document}
